% !PRINTED 431
\weberpage{431}

\S\ 119. \gesperrt{Die Diskriminante der Invariantengleichung.}

Wir stellen uns jetzt die Frage, wann eine Invariantengleichung
\weq{1}{F_p(v,\,u) = 0,}
worin $p$ eine Primzahl ist, für $u = j(\omega)$ zwei oder mehrere gleiche Wurzeln $v$ hat, oder wann unter den $p+1$ Größen
\weq{2}{j(p\,\omega),\quad j\!\left(\frac{\omega+c}{p}\right),\quad c = 0,\,1,\,\ldots,\,p-1}
zwei oder mehrere einander gleiche vorkommen. Dies findet dann und nur dann statt, wenn unter den $p+1$ Größen
\weq{3}{p\,\omega,\quad \frac{\omega+c}{p}}
zwei äquivalente vorkommen. Es muß also jedenfalls $\omega$ einer ganzzahligen quadratischen Gleichung mit negativer Diskriminante $D=-m$ genügen, die wir in der Form annehmen:
\weq{4}{A\,\omega^2 + B\,\omega + C = 0,\quad 4\,A\,C - B^2 = m > 0,}
worin $A,\,B,\,C$ ohne gemeinsamen Teiler sind.

Ersetzen wir $\omega$ durch eine äquivalente Zahl, also (4) durch eine äquivalente Gleichung, so werden die Größen (2) nur untereinander vertauscht. Die Frage nach der \gesperrt{Anzahl} der gleichen Wurzeln von (1) wird also davon nicht berührt, und wir können daher annehmen, daß $A$ durch $p$ nicht teilbar sei. Ist nun zunächst $p\,\omega$ äquivalent mit $\dfrac{\omega+c}{p}$, so ist:
\weq{5}{\frac{\omega+c}{p} = \frac{\gamma+\delta\,p\,\omega}{\alpha+\beta\,p\,\omega},\quad \alpha\,\delta-\beta\,\gamma = 1.}
Schreiben wir diese Gleichung so:
\weq{6}{\beta\,p\,\omega^2 + (\alpha+\beta\,c\,p-\delta\,p^2)\,\omega + c\,\alpha-\gamma\,p = 0,}
so folgt durch Vergleichung mit (4) (da $A$ durch $p$ unteilbar sein sollte), daß $\alpha$ durch $p$ teilbar, also $=p\,\alpha'$ sein muß, und daß eine ganze Zahl $x$ existiert, die den Bedingungen genügt:
\weq{7}{\beta = A\,x,\quad \alpha'+\beta\,c-\delta\,p = B\,x,\quad c\,\alpha'-\gamma = C\,x.}
Setzen wir noch
\weq{8}{\alpha'-\beta\,c+\delta\,p = y,}
so folgt:
\weq{9}{\begin{aligned}
2\,\alpha' &= B\,x+y,\\
2\,(\beta\,c-\delta\,p) &= B\,x-y.
\end{aligned}}
%CONT>>
